% ===================================================================
%  TABELO N-RO 1 — La lernejo. — La liceo. — La klasĉambro.
%
%  Ceci n'est PAS une transcription : c'est une traduction, faite en
%  2026, en esperanto d'aujourd'hui. D'ou trois differences avec
%  texto/io et texto/fr, et trois seulement :
%
%   * pas de \nl ni de \cc — la traduction ne suit aucune ligne
%     imprimee, et les lignes de ce fichier ne sont que du confort ;
%   * pas de \begin{VUpage} — elle n'a ni feuillet ni folio, et la
%     page de lecture ne lui en affiche pas ;
%   * le gras se pose sur le TERME seul, ponctuation dehors.
%
%  Tout le reste est identique, et doit l'etre : les cles « %%K » sont
%  celles de texto/io, macro pour macro pour l'apparat, et les renvois
%  \textsuperscript{(n)} visent les MEMES objets de la planche — ce
%  sont eux qui ouvrent les gros plans.
%
%  LA PREMIERE DES LANGUES QUE LES IDISTES PARLENT DEJA, et la plus
%  proche : l'ido est ne de l'esperanto en 1907. On traduit donc SUR
%  L'IDO directement, non sur une colonne intermediaire, et l'ordre
%  des renvois se garde presque sans effort — les deux langues rangent
%  leurs mots de la meme facon.
%
%  ET C'EST JUSTEMENT POURQUOI IL FAUT SE MEFIER. Un mot d'ido n'est
%  pas un mot d'esperanto parce qu'il lui ressemble : « docisto » se
%  dit ici « instruisto », « lernanto » reste « lernanto », mais
%  « docochambro » devient « klasĉambro », « skolo » devient
%  « lernejo », « stulo » devient « seĝo », « plaso » devient
%  « loko ». La ressemblance se verifie mot par mot, elle ne se
%  suppose pas.
%
%  L'ACCUSATIF ET L'ACCORD sont de l'esperanto et non de l'ido :
%  « la instruisto legas la libron », « bonaj lernantoj ». Ils ne
%  deplacent aucun renvoi — le complement reste apres son verbe — mais
%  ils changent la fin de presque chaque mot.
%
%  LES NOMS DE PERSONNE PRENNENT LEUR FORME ESPERANTISTE, comme dans
%  les treize autres colonnes : Ioannes est Johano, Iakobus est
%  Jakobo, Karolus est Karlo. Seuls le chien Medoro et l'ane Aliborono
%  gardent leur nom. Les noms latins de la note (*) ne bougent pas :
%  c'est d'eux que la note parle.
% ===================================================================

%%K t01-apar-0 apar
\VUsaut{2.0mm}
\VUcentre{12.6pt}{120}{EKSPLIKA LIBRETO}
\VUsaut{3.2mm}
\VUcentre{8.4pt}{160}{DE}
\VUsaut{2.6mm}
\VUtitre{15.6pt}{0}{la helpaj tabeloj Delmas}
\VUsaut{3.4mm}
\VUornamento{9pt}
\VUsaut{5.0mm}
\VUcentre{13.2pt}{90}{UNUA SERIO}
\VUsaut{3.0mm}
\VUfilet{16mm}
\VUsaut{5.6mm}

%%K t01-apar-1 apar
\VUtitre{15.0pt}{60}{TABELO N\textsuperscript{o} 1}
\VUsaut{1.4mm}
\VUtitre{11.4pt}{0}{La lernejo, la liceo, la klasĉambro.}
\VUsaut{2.2mm}
\VUfilet{20mm}
\VUsaut{6.4mm}

%%K t01-tit-1 sub
\VUpk{10.2pt}{40}{Ĝenerala priskribo.}
\VUsaut{3.0mm}

%%K t01-01-1 p
1. --- Ĉi tiu tabelo prezentas \VUgras{klasĉambron} en \VUgras{liceo}. En
tiu klasĉambro estas ok \VUgras{tabloj} \textsuperscript{(18)} kaj ok
\VUgras{benkoj} \textsuperscript{(32)}. Sur ĉiu benko sidas tri lernantoj. La
\VUgras{instruisto} \textsuperscript{(7)} sidas sur \VUgras{seĝo}
\textsuperscript{(8)} en sia \VUgras{katedro} \textsuperscript{(6)}.

%%K t01-02-1 p
2. --- Maldekstre de la katedro troviĝas \VUgras{vitrita} \VUgras{ŝranko}
\textsuperscript{(15)}. Dekstre estas la \VUgras{nigra tabulo}
\textsuperscript{(1)} kun la \VUgras{viŝilo} \textsuperscript{(2)} kaj la
\VUgras{kreto} \textsuperscript{(3)}.

%%K t01-02-2 p
Super la katedro pendas \VUgras{muraj tabeloj} \textsuperscript{(9, 11, 12)}
kaj \VUgras{mapo} \textsuperscript{(10)}.

%%K t01-03-1 p
3. --- Ne ĉiuj lernantoj sidas en sia \VUgras{loko} \textsuperscript{(17)}.
Johano \textsuperscript{(4)} staras ĉe la tabulo, li tenas la viŝilon kaj
forviŝas la \VUgras{geometriajn figurojn}.

%%K t01-03-2 p
Petro estas proksime de la katedro, li kolektas la \VUgras{skribotaskojn}
\textsuperscript{(5)} kaj alportas ilin al la instruisto.

%%K t01-04-1 p
4. --- Ĉi tiu instruisto estas la instruisto pri la \VUgras{vivantaj lingvoj}.
Li instruas al la lernantoj paroli, legi kaj skribi fremdajn lingvojn
\VUgras{(la germanan, la anglan, la hispanan, la italan, la rusan} aŭ
\VUgras{la francan)}.

%%K t01-05-1 p
5. --- La klasĉambro estas tre vasta kaj tre hela. En la fundo estas larĝa
\VUgras{fenestro} kun du \VUgras{luketoj} \textsuperscript{(78)} kaj
\VUgras{vitrita pordo} por aerumi kaj lumigi ĝin.

%%K t01-05-2 p
Ĉi tiu klasĉambro ne estas malvarma. En la mezo estas por varmigi ĝin tre bona
\VUgras{forno} \textsuperscript{(46)} kun sia \VUgras{tubo}
\textsuperscript{(45)}.

%%K t01-05-3 p
Ĉe la muro estas fiksitaj \VUgras{vesthokoj} \textsuperscript{(58)}, de ili
pendas la ĉapoj kaj la manteloj de la lernejanoj.

%%K t01-tit-2 sub
\VUsaut{5.2mm}
\VUpk{10.2pt}{40}{La ludokorto}
\VUsaut{3.6mm}

%%K t01-06-1 p
6. --- La \VUgras{horloĝo} \textsuperscript{(63)} montras naŭ horojn kaj kvin
minutojn.

%%K t01-06-2 p
La \VUgras{paŭzo} \textsuperscript{(76)} ĵus finiĝis kaj la leciono
rekomenciĝas. La paŭzejo troviĝas en la \VUgras{korto} \textsuperscript{(75)}.
Ni vidas kelkajn lernantojn kiuj ludas. \VUgras{Submajstro}
\textsuperscript{(74)} kontrolas ilin.

%%K t01-07-1 p
7. --- En la korto ni krome vidas diversajn personojn, nome la
\VUgras{inspektiston} \textsuperscript{(73)}, la lernejestron
\VUgras{(liceestron)} \textsuperscript{(71)} kaj la \VUgras{cenzoron}
\textsuperscript{(72)}.

%%K t01-07-2 p
La inspektisto inspektas la lernejojn, la liceestro direktas la liceon, la
cenzoro kontrolas la studojn.

%%K t01-07-3 p
La kvara persono estas la \VUgras{pordisto} \textsuperscript{(77)}. Li portas
sub la brako registron por enskribi la forestantojn.

%%K t01-08-1 p
8. --- En la fundo de la korto estas la \VUgras{gimnastikaj aparatoj}
\textsuperscript{(70)} kaj la laborejo por la \VUgras{manlaboroj}
\textsuperscript{(69)}.

%%K t01-08-2 p
Dekstre de la tabelo ni ekvidas la \VUgras{klasĉambrojn} pri \VUgras{muziko}
kaj \VUgras{desegno}.

%%K t01-noto-1 noto
\VUnotes{143.2mm}{%
(*) Por la \textit{baptonomoj} oni konsilas ankaŭ transskribi ilin laŭ la
latina formo. Tio estas la sola rimedo eskapi sennombrajn diverĝojn. Cetere
kiel elekti inter \textit{Giovanni, Jean,} \textit{Johann, Jan, Hans, John,
Ivan,} k. c.? Ĉu ni kreos specialan vortaron por la baptonomoj? Ni ĉerpu el la
latina ilian nominativon kaj diru : \VUgras{Ioannes, Ioanna; Iulius, Iulia;
Iakobus; Andreas; Lukas.} (Kompleta Gram.).%
}

%%K t01-tit-3 sub
\VUpk{10.2pt}{40}{Detala priskribo.}
\VUsaut{2.4mm}

%%K t01-tit-4 sub
\VUpk{10.2pt}{40}{La bonaj lernantoj.}
\VUsaut{3.0mm}

%%K t01-09-1 p
9. --- La lernantoj de la unua benko dekstre de la katedro estas
\VUgras{Henriko} (*), \VUgras{Stefano} kaj \VUgras{Karlo}.

%%K t01-09-2 p
Henriko estas tre atentema kaj laborema, li aŭskultas la instruiston.

Sur sia tablo li havas \VUgras{kajeron} \textsuperscript{(28)},
\VUgras{libron} \textsuperscript{(29)}, \VUgras{vortaron}
\textsuperscript{(30)} kaj \VUgras{plumujon} \textsuperscript{(31)}. Lia libro
havas multajn \VUgras{paĝojn}, ĉiu ĉapitro enhavas plurajn
\VUgras{paragrafojn} kaj ĉiu \VUgras{linio} multajn \VUgras{vortojn}.

%%K t01-09-4 p
Lia najbaro Stefano \textsuperscript{(24)} estas fervora. Li notas sur sia
kajero la lecionon por la venonta fojo. Li havas belan \VUgras{skribmanieron}.
Lia libro estas fermita. Ni vidas nur la \VUgras{kovrilon}
\textsuperscript{(27)}. Apud li troviĝas lia \VUgras{cirkelo}
\textsuperscript{(26)}.

%%K t01-09-5 p
Karlo \textsuperscript{(23)} tenas per la mano sian libron, li malfermas ĝin
por relegi sian lecionon. Li havas, kiel ĉiu lernanto, \VUgras{inkujon}
\textsuperscript{(25)} antaŭ si. Li estas obeema.

%%K t01-10-1 p
10. --- Ĉe la apuda tablo estas \VUgras{Ludoviko, Antono} kaj \VUgras{Paŭlo}.
Ludoviko recitas sian \VUgras{lecionon} \textsuperscript{(22)} kaj respondas al
la \VUgras{demandoj} de la instruisto. Antono \textsuperscript{(21)} estas tre
neatentema, eĉ kelkfoje la instruisto punas lin. Paŭlo \textsuperscript{(20)}
estas bona kamarado, afabla kaj komplezema. Li estas tre atentema.

%%K t01-tit-5 sub
\VUsaut{4.2mm}
\VUpk{10.2pt}{40}{La malbonaj lernantoj.}
\VUsaut{3.2mm}

%%K t01-11-1 p
11. --- Malantaŭ Henriko troviĝas \VUgras{Georgo} \textsuperscript{(38)}. Li ne
estas malbona knabo, sed li estas \VUgras{babilema}, neatentema kaj petolema.
Lia \VUgras{frivoleco} kaj \VUgras{malobeemo} kaŭzas \VUgras{malordon} dum la
leciono kaj ofte li meritas severan \VUgras{averton}. Lia \VUgras{krudkajero}
\textsuperscript{(35)} estas malpura, li \VUgras{inkmakulas} ĝin per sia
\VUgras{plumo} \textsuperscript{(34)}, pro tio li ĉiam uzas sian
\VUgras{skrapgumon} \textsuperscript{(36)}; lia \VUgras{kajerujo}
\textsuperscript{(33)} estas ŝirita, ĉar li estas tre neglektema. Kial li
kunportas sian \VUgras{krajontenilon} \textsuperscript{(37)} en la klasĉambron
de la vivantaj lingvoj?

%%K t01-11-2 p
Lia najbaro Edmundo \textsuperscript{(39)} ne amas lin. Li estas vere iom
maldiligenta, sed li estas silentema kaj trankvila. Li ne babilas; sed,
anstataŭ aŭskulti, li preferas legi la \VUgras{rakontojn} de sia
\VUgras{legolibro}.

%%K t01-11-3 p
La tria lernanto de ĉi tiu benko estas \VUgras{Ludoviko}
\textsuperscript{(40)}. Li estas diligenta. Li skribas sur blanka
\VUgras{paperfolio}. Antaŭ si li metis sian \VUgras{sorbpaperon}
\textsuperscript{(41)}.

%%K t01-12-1 p
12. --- La lernantoj de la dua benko, maldekstre de la instruisto, estas tre
junaj. La unua, la malgranda \VUgras{Emilio}, ankoraŭ ne scias tre bone
skribi, li kunportas sian \VUgras{ardezon} \textsuperscript{(42)}, sian
\VUgras{ardezan krajonon} kaj sian \VUgras{spongon} \textsuperscript{(43)}.

%%K t01-12-2 p
Apud li estas \VUgras{Aŭgusto} kaj \VUgras{Bernardo} \textsuperscript{(44)}.
Ĉi tiu bone laboras, sed lia \VUgras{vestiteco} estas tre neglektita.

%%K t01-13-1 p
13. --- Malantaŭ ili troviĝas tri \VUgras{maldiligentuloj}. \VUgras{Alberto}
ne lernas siajn lecionojn. \VUgras{Jakobo} \textsuperscript{(47)} dormas
anstataŭ aŭskulti. Lia \VUgras{maldiligenteco} kaŭzas al li
\VUgras{riproĉojn} kaj eĉ \VUgras{punojn}. Fine \VUgras{Kristiano}
\textsuperscript{(48)} estas tro neserioza. Li malbone \VUgras{kondutas} kaj
neniel penas por korekti sin.

%%K t01-14-1 p
14. --- Liaj najbaroj, kontraste, estas bonkondutaj. \VUgras{Ernesto}
\textsuperscript{(51)} ŝatas la ordon kaj la regulecon, li ĉiam uzas sian
\VUgras{linialon} \textsuperscript{(50)} kaj sian \VUgras{krajonon}
\textsuperscript{(49)}. Li estas \VUgras{eksterulo}.

%%K t01-14-2 p
\VUgras{Francisko} \textsuperscript{(52)} estas \VUgras{internulo}, li estas
vestita per uniformo, manĝas kaj kuŝas en la liceo. Li estas bona
\VUgras{kamarado}.

%%K t01-14-3 p
Maŭrico akrigas sian \VUgras{krajonon} per sia \VUgras{tranĉileto}
\textsuperscript{(56)}, li estas tre zorgema.

%%K t01-14-4 p
Li kunportas ĉiujn siajn librojn, sian \VUgras{gramatikon}
\textsuperscript{(55)}, sian \VUgras{notlibreton} \textsuperscript{(54)}, kaj
eĉ \VUgras{skribkusenon} \textsuperscript{(53)}. Lia \VUgras{lernosako}
\textsuperscript{(57)} estas sur la planko malantaŭ li.

%%K t01-14-5 p
La du tabloj en la fundo de la klasĉambro estas okupataj de \VUgras{bonaj
lernantoj} \textsuperscript{(19)}.

%%K t01-tit-6 sub
\VUsaut{4.6mm}
\VUpk{10.2pt}{40}{Desegno kaj muziko.}
\VUsaut{3.4mm}

%%K t01-15-1 p
15. --- En la \VUgras{desegnoĉambro} estas graciaj \VUgras{modeloj} pendigitaj
ĉe la muro kaj \VUgras{lampo} \textsuperscript{(62)} kun \VUgras{ŝirmilo},
suspendita ĉe la plafono. La \VUgras{instruisto} \textsuperscript{(60)} estas
tre pacienca, li korektas la \VUgras{desegnaĵojn} \textsuperscript{(61)} de la
lernantoj kaj instruas al ili desegni \VUgras{buston} \textsuperscript{(59)}
laŭ la naturo.

%%K t01-16-1 p
16. --- La \VUgras{muzikejo} ŝajnas tre interesa. En ĝi oni lernas legi la
\VUgras{muzikon}, solfeĝi la \VUgras{gamnotojn} \textsuperscript{(67)}
skribitajn sur la \VUgras{liniaro} (do, re, mi, fa, sol, la, si\,=\,C. D. E.
F. G. A. B.), koni la \VUgras{kleojn} kaj kanti ĉarmajn \VUgras{melodiojn}.
Oni metas la muzikaĵojn sur la \VUgras{pupitron} \textsuperscript{(65)}. Unu el
la lernantoj \VUgras{ludas pianon} \textsuperscript{(64)} kaj la
\VUgras{muzikinstruisto} \textsuperscript{(66)} \VUgras{taktas}
\textsuperscript{(68)}.

%%K t01-17-1 p
17. --- Ni ekvidas angulon de la \VUgras{laborejo} por \VUgras{manlaboroj}
\textsuperscript{(69)}, kie la knaboj rajtas lerni raboti la lignon kaj fajli
la feron. Tio ne ekzistas en ĉiuj liceoj.

%%K t01-tit-7 sub
\VUsaut{5.0mm}
\VUpk{10.2pt}{40}{La instruejo pri la sciencoj.}
\VUsaut{3.6mm}

%%K t01-18-1 p
18. --- La instruisto pri matematiko instruas al la lernantoj
\VUgras{geometrion, aritmetikon, kalkulon} kaj la \VUgras{metran sistemon}. Ni
vidas sur la tabelo la ĉefajn figurojn de la geometrio :
\VUgras{cirklon} \textsuperscript{(a)}, \VUgras{kvadraton}
\textsuperscript{(b)}, \VUgras{triangulon} \textsuperscript{(c)},
\VUgras{linion} \textsuperscript{(d)}.

%%K t01-18-2 p
Ni vidas ankaŭ \VUgras{kalkulon}; ĝi estas nek \VUgras{adicio}, nek
\VUgras{subtraho}, nek \VUgras{multipliko} : ĝi estas \VUgras{divido}
\textsuperscript{(e)}.

%%K t01-19-1 p
19. --- Sur la tabelo de la metra sistemo ni vidas : \VUgras{metron}
\textsuperscript{(a)}, \VUgras{pesilon} \textsuperscript{(b)} kaj ĝiajn
\VUgras{pezilojn}, \VUgras{litron} \textsuperscript{(e)},
\VUgras{dekalitron} \textsuperscript{(f)}, \VUgras{decilitron} kaj
\VUgras{kubametron} \textsuperscript{(d)}.

%%K t01-19-2 p
La lernantoj studas ankaŭ la \VUgras{naturajn sciencojn}
\textsuperscript{(11)} --- zoologion \textsuperscript{(a)}, botanikon
\textsuperscript{(b)}, geologion \textsuperscript{(c)} --- kaj la
\VUgras{historion} \textsuperscript{(12)}.

%%K t01-19-3 p
En la ŝranko estas la aparatoj de \VUgras{fiziko} \textsuperscript{(15)} kaj de
\VUgras{ĥemio} \textsuperscript{(16)}.

%%K t01-19-4 p
Super la ŝranko estas : \VUgras{sfero} \textsuperscript{(14)}, \VUgras{konuso}
kaj \VUgras{terglobo} \textsuperscript{(13)}.

%%K t01-19-5 p
Super la tabeloj pendas \VUgras{mapo} kiu montras la kvin partojn de la mondo :
\VUgras{Eŭropon} \textsuperscript{(g)}, \VUgras{Azion} \textsuperscript{(e)},
\VUgras{Afrikon} \textsuperscript{(f)}, \VUgras{Amerikon}
\textsuperscript{(ab)} kaj \VUgras{Oceanion} \textsuperscript{(h)}, kaj la du
grandajn oceanojn, la \VUgras{Pacifikan} \textsuperscript{(c)} kaj la
\VUgras{Atlantikan} \textsuperscript{(d)}.
\VUsaut{6.0mm}
\VUfilet{22mm}
